% ============================================================
% 图论算法示例 —— Dijkstra、BFS、并查集
% 代码来自 sy4-4 天梯地图 / sy4-8 周游世界 / sy4-1 并查集
% 编译方式: xelatex 图论算法示例.tex
% ============================================================
\documentclass[12pt,a4paper]{article}

\usepackage{geometry}
\geometry{left=2.0cm,right=2.0cm,top=2.0cm,bottom=2.0cm}

\usepackage{ctex}
\usepackage{listings}
\usepackage{xcolor}
\usepackage{hyperref}

% ===== 代码样式（与 study/CS/hyper/ 一致）=====
\lstset{
    language=C++,
    basicstyle=\small\ttfamily,
    keywordstyle=\bfseries\color{blue!75!black},
    commentstyle=\color{green!60!black},
    stringstyle=\color{red!70!black},
    numbers=left,
    numberstyle=\tiny\color{gray},
    numbersep=8pt,
    frame=none,
    breaklines=true,
    tabsize=4,
    showstringspaces=false,
    extendedchars=true,
    inputencoding=utf8,
    aboveskip=6pt,
    belowskip=6pt
}

\usepackage{titlesec}
\titleformat{\section}{\Large\bfseries}{\thesection}{0.5em}{}
\titleformat{\subsection}{\large\bfseries}{\thesubsection}{0.5em}{}

\begin{document}

\begin{center}
    {\Huge\bfseries 图论算法示例}\\[0.3cm]
    {\large 数据结构考试复习 --- Dijkstra / BFS / 并查集}\\[0.2cm]
    {2026 年 6 月 18 日}
\end{center}
\vspace{0.3cm}

\tableofcontents
\newpage

% ============================================================
% 1  Dijkstra 算法
% ============================================================
\section{Dijkstra 算法（天梯地图）}

\subsection{图的定义}

\begin{lstlisting}
#include <iostream>
#include <vector>
#include <climits>   // 用 INT_MAX 定义无穷大
#include <algorithm> // 用 reverse 还原路径
using namespace std;

// 用一个足够大的数表示无穷大
// 使用 INT_MAX/2 而不是 INT_MAX，防止 dist[u] + w 时溢出
const int INF = INT_MAX / 2;

// 边的结构体，表示图中的一条有向边
struct Edge {
    int to;          // 目标节点编号
    int length;      // 道路长度（距离），单位千米
    int time;        // 通过该道路所需时间，单位分钟
};

// 节点的结构体，表示图中的一个顶点
struct Node {
    vector<Edge> edges;  // 从该节点出发的所有有向边
};
\end{lstlisting}

\subsection{Dijkstra 核心框架（双目标优化）}

\begin{lstlisting}
/**
 * Dijkstra 求最快到达路线
 * 主目标：时间最短
 * 次目标：时间相同时选距离更短的
 *
 * @param n     节点总数
 * @param start 起点编号
 * @param end   终点编号
 * @param graph 图的邻接表（vector<Node>）
 * @return pair<路径向量, 总时间>
 */
pair<vector<int>, int> dijkstraTime(int n, int start, int end,
                                    const vector<Node>& graph) {
    // distTime[i]：起点到节点 i 的最短时间
    vector<int> distTime(n, INF);
    // distLen[i]：在最短时间条件下到节点 i 的最短距离（打破平局用）
    vector<int> distLen(n, INF);
    // parent[i]：节点 i 的前驱节点，用于最后还原完整路径
    vector<int> parent(n, -1);
    // visited[i]：标记节点 i 是否已确定最短路径
    vector<bool> visited(n, false);

    // 起点到自身的距离和时间均为 0
    distTime[start] = 0;
    distLen[start] = 0;

    // 主循环：每次选出一个未访问且时间最小的节点
    for (int i = 0; i < n; i++) {
        // 在未访问节点中找 distTime 最小的
        int u = -1;
        int minTime = INF;
        for (int j = 0; j < n; j++) {
            if (!visited[j] && distTime[j] < minTime) {
                minTime = distTime[j];
                u = j;
            }
        }
        // 如果所有剩余节点都不可达，提前结束
        if (u == -1) break;
        visited[u] = true;
        // 已到达终点，可以提前退出
        if (u == end) break;

        // 松弛操作：用节点 u 更新其所有邻居的距离
        for (const Edge& e : graph[u].edges) {
            int v = e.to;
            if (!visited[v]) {
                // 情况1：找到更短的时间 → 直接更新
                if (distTime[u] + e.time < distTime[v]) {
                    distTime[v] = distTime[u] + e.time;
                    distLen[v] = distLen[u] + e.length;
                    parent[v] = u;  // 记录前驱
                }
                // 情况2：时间相同但距离更短 → 同样更新
                else if (distTime[u] + e.time == distTime[v]
                        && distLen[u] + e.length < distLen[v]) {
                    distLen[v] = distLen[u] + e.length;
                    parent[v] = u;
                }
            }
        }
    }

    // 返回路径和总时间
    return {getPath(parent, start, end), distTime[end]};
}
\end{lstlisting}

\subsection{路径还原}

\begin{lstlisting}
/**
 * 从 parent 数组还原完整路径
 * parent[i] 记录节点 i 的前驱节点，从终点向前追溯到起点
 */
vector<int> getPath(const vector<int>& parent, int start, int end) {
    vector<int> path;
    int cur = end;
    // 从终点开始，通过 parent 向前追溯直到起点
    // parent[start] == -1（起点没有前驱）
    while (cur != -1) {
        path.push_back(cur);
        cur = parent[cur];
    }
    // 上面得到的是从终点到起点的逆序，反转后得到正确顺序
    reverse(path.begin(), path.end());
    return path;
}

/**
 * 格式化输出路径：start => v1 => v2 => ... => end
 */
void printPath(const vector<int>& path, int start, int end) {
    cout << start;  // 先输出起点
    // 依次输出路径上的后续节点，用" => "连接
    for (int i = 1; i < path.size(); i++)
        cout << " => " << path[i];
}
\end{lstlisting}

\subsection{主函数与图的构建}

\begin{lstlisting}
int main() {
    int n, m;  // n: 节点数, m: 边数
    cin >> n >> m;

    // 创建邻接表表示的图，共 n 个节点
    vector<Node> graph(n);

    // 读入 m 条边
    for (int i = 0; i < m; i++) {
        int v1, v2, oneway, length, time;
        cin >> v1 >> v2 >> oneway >> length >> time;

        // 正向边 v1 -> v2
        Edge e1 = {v2, length, time};
        graph[v1].edges.push_back(e1);

        // oneway == 0 表示双向路，添加反向边 v2 -> v1
        if (oneway == 0) {
            Edge e2 = {v1, length, time};
            graph[v2].edges.push_back(e2);
        }
        // oneway == 1 表示单行道，不添加反向边
    }

    int start, end;  // 查询的起点和终点
    cin >> start >> end;

    // 调用 Dijkstra 计算最快路线
    auto [timePath, totalTime] =
        dijkstraTime(n, start, end, graph);

    // 输出结果
    cout << "Time = " << totalTime << ": ";
    printPath(timePath, start, end);
    cout << endl;

    return 0;
}
\end{lstlisting}

\newpage

% ============================================================
% 2  BFS 广度优先搜索
% ============================================================
\section{BFS 广度优先搜索（周游世界）}

\subsection{图的定义}

\begin{lstlisting}
#include <iostream>
#include <vector>
#include <queue>     // BFS 需要队列
#include <algorithm> // 用 reverse 还原路径
using namespace std;

// 边的结构体
struct Edge {
    int to;          // 目标站点编号
    int company;     // 这条边属于哪家公司（用于计算换乘）
};
\end{lstlisting}

\subsection{BFS 求最少经停站（无权图最短路径）}

\begin{lstlisting}
/**
 * BFS 求无权图中从 start 到各点的最少边数（经停站数）
 * 核心思想：BFS 按层遍历，第一次访问到某节点时即为最短距离
 *
 * @param start 起点
 * @param n     节点总数
 * @param graph 邻接表
 * @return dist 数组，dist[i] = 从 start 到 i 的最少边数，-1 表示不可达
 */
vector<int> bfsShortestPath(int start, int n,
                            const vector<vector<Edge>>& graph) {
    // dist[i] = -1 表示尚未访问（也即不可达）
    vector<int> dist(n, -1);
    queue<int> q;

    // 起点到自身距离为 0
    dist[start] = 0;
    q.push(start);

    // BFS 主循环：队列非空时持续处理
    while (!q.empty()) {
        int u = q.front();  // 取队首节点
        q.pop();            // 队首出队

        // 遍历节点 u 的所有邻居
        for (const Edge& e : graph[u]) {
            int v = e.to;
            // 如果邻居 v 尚未访问，说明这是到 v 的最短路
            if (dist[v] == -1) {
                dist[v] = dist[u] + 1;  // 距离 = 前驱距离 + 1
                q.push(v);              // v 入队，后续继续扩展
            }
        }
    }
    return dist;
    // 查询结果：dist[end] 即为 start 到 end 的最少经停站数
    // 若 dist[end] == -1，表示不可达
}
\end{lstlisting}

\subsection{BFS 多条件（经停站最少 + 换乘最少）}

\begin{lstlisting}
/**
 * PreNode 结构体：记录 BFS 路径上的前驱信息
 * 用于在找到目标后还原完整路径，以及判断换乘
 */
struct PreNode {
    int from;          // 前驱节点编号
    int company;       // 从前驱节点到当前节点经过的公司
    int transfers;     // 到当前节点为止的累计换乘次数
};

/**
 * BFS 同时求最少经停站和最少换乘
 * 优先保证经停站最少，经停站相同时取换乘最少的路径
 */
void bfsWithTransfer(int start, int target,
                     const vector<vector<Edge>>& graph) {
    int n = graph.size();

    // dist[i]：从 start 到 i 的最少经停站数
    vector<int> dist(n, -1);
    // transfers[i]：在经停站最少的前提下，到 i 的最少换乘次数
    vector<int> transfers(n, 1e9);
    // pre[i]：在最优路径上 i 的前驱信息
    vector<PreNode> pre(n, {-1, -1, 0});

    queue<int> q;
    dist[start] = 0;
    transfers[start] = 0;
    q.push(start);

    while (!q.empty()) {
        int u = q.front();
        q.pop();

        // 遍历 u 的所有邻居
        for (const Edge& e : graph[u]) {
            int v = e.to;
            int comp = e.company;

            // 计算到达 v 的新换乘次数
            int newTransfers = transfers[u];
            if (u != start) {
                // 如果当前边的公司与前一站的公司不同，说明换乘了一次
                if (pre[u].company != comp)
                    newTransfers++;
            }
            // 注意：如果 u 是起点，第一段线路不算换乘

            // 判断是否更新 v 的最优值：
            // 条件1：v 尚未被访问过
            // 条件2：找到了更少的经停站数
            // 条件3：经停站数相同但换乘更少
            if (dist[v] == -1 ||
                dist[u] + 1 < dist[v] ||
                (dist[u] + 1 == dist[v] && newTransfers < transfers[v])) {

                dist[v] = dist[u] + 1;
                transfers[v] = newTransfers;
                pre[v] = {u, comp, transfers[v]};
                q.push(v);
            }
        }
    }

    // 输出结果
    if (dist[target] == -1) {
        cout << "Sorry, no line is available." << endl;
    } else {
        // 最少经停站数（不含起点，dist 统计的是边数）
        cout << dist[target] << endl;

        // 回溯还原完整路径
        vector<int> path;
        int cur = target;
        while (cur != -1) {
            path.push_back(cur);
            cur = pre[cur].from;
        }
        reverse(path.begin(), path.end());

        // 输出路径全长（含起点和终点）
        cout << "路径长度: " << path.size() << endl;
    }
}
\end{lstlisting}

\subsection{BFS 求连通分量}

\begin{lstlisting}
/**
 * 从一个起点开始 BFS，标记所有可达节点
 *
 * @param start     起点
 * @param adj       邻接表（无权图）
 * @param visited   访问标记数组（传引用，会修改）
 * @param traversal 存储遍历顺序（传引用）
 */
void bfs(int start, const vector<vector<int>>& adj,
         vector<bool>& visited, vector<int>& traversal) {
    queue<int> q;
    q.push(start);
    visited[start] = true;

    while (!q.empty()) {
        int u = q.front();
        q.pop();
        traversal.push_back(u);  // 记录遍历顺序

        // 遍历所有邻接点
        for (int v : adj[u]) {
            if (!visited[v]) {
                visited[v] = true;
                q.push(v);
            }
        }
    }
}

/**
 * 对非连通图找所有连通分量
 * 遍历每个节点，如果未访问就以它为起点做一次 BFS
 */
void findAllComponents(int n, const vector<vector<int>>& adj) {
    vector<bool> visited(n, false);
    vector<int> traversal;  // 按 BFS 顺序存储所有节点
    int components = 0;     // 连通分量计数

    for (int i = 0; i < n; i++) {
        if (!visited[i]) {
            bfs(i, adj, visited, traversal);
            components++;  // 每次 BFS 发现一个新的连通分量
        }
    }

    cout << "连通分量数: " << components << endl;
    // traversal 中存储了 BFS 遍历序列
}
\end{lstlisting}

\newpage

% ============================================================
% 3  并查集（Union-Find）
% ============================================================
\section{并查集（Union-Find / Disjoint Set Union）}

\subsection{结构体实现（路径压缩 + 按秩合并）}

\begin{lstlisting}
#include <iostream>
#include <vector>
#include <algorithm>
using namespace std;

const int MAXN = 10000;

/**
 * UnionFind 并查集结构体
 * 支持三种操作：
 *   find(x)  —— 查找 x 所在集合的根节点
 *   unite(x, y) —— 合并 x 和 y 所在的集合
 *   same(x, y) —— 判断 x 和 y 是否属于同一集合
 */
struct UnionFind {
    // parent[i]：节点 i 的父节点，根节点的父节点指向自身
    vector<int> parent;
    // rank[i]：以 i 为根的树的高度（近似），用于按秩合并优化
    vector<int> rank;

    /**
     * 构造函数：初始化 n 个独立元素
     * 每个元素的父节点指向自己，即每个元素自成一族
     */
    UnionFind(int n) {
        parent.resize(n);
        rank.resize(n, 0);
        for (int i = 0; i < n; i++)
            parent[i] = i;
    }

    /**
     * 查找元素 x 所在集合的根节点
     * 同时进行路径压缩：将路径上所有节点直接指向根节点
     * 使得后续查找更快（近似 O(1)）
     */
    int find(int x) {
        if (parent[x] != x)
            parent[x] = find(parent[x]);  // 递归路径压缩
        return parent[x];
    }

    /**
     * 合并 x 和 y 所在的集合
     * 按秩合并：将较矮的树挂在较高的树下，防止树退化成链表
     */
    void unite(int x, int y) {
        int rx = find(x), ry = find(y);
        if (rx == ry) return;  // 已经在同一集合，无需合并

        // 按秩合并：rank 小的挂到 rank 大的下面
        if (rank[rx] < rank[ry])
            parent[rx] = ry;      // rx 的根指向 ry
        else if (rank[rx] > rank[ry])
            parent[ry] = rx;      // ry 的根指向 rx
        else {
            // rank 相等时，合并后 rank 增加 1
            parent[ry] = rx;
            rank[rx]++;
        }
    }

    /**
     * 判断 x 和 y 是否属于同一集合
     * 根节点相同则在同一集合
     */
    bool same(int x, int y) {
        return find(x) == find(y);
    }
};
\end{lstlisting}

\subsection{简化数组版（竞赛常用）}

\begin{lstlisting}
// 全局父节点数组，parent[i] 表示节点 i 的父节点
int parent[MAXN];

/**
 * 初始化：每个元素自成一族
 * 即每个元素的父节点指向自身
 */
void init(int n) {
    for (int i = 0; i < n; i++)
        parent[i] = i;
}

/**
 * 查找根节点（路径压缩）
 * 递归查找的同时压缩路径，使树保持扁平
 */
int find(int x) {
    if (parent[x] != x)
        parent[x] = find(parent[x]);
    return parent[x];
}

/**
 * 合并两个元素所在的集合
 * 将 x 的根直接指向 y 的根（或反过来）
 */
void unite(int x, int y) {
    int rx = find(x), ry = find(y);
    if (rx != ry)
        parent[ry] = rx;  // 将 ry 的根设为 rx
}
\end{lstlisting}

\subsection{应用：统计连通分量}

\begin{lstlisting}
/**
 * 统计当前图中的连通分量数
 * 原理：根节点（parent[i] == i）的个数 == 连通分量数
 */
int countComponents(int n) {
    int cnt = 0;
    for (int i = 0; i < n; i++)
        if (find(i) == i)  // 自己是根节点 → 一个分量
            cnt++;
    return cnt;
}
\end{lstlisting}

\subsection{应用：按根节点分组}

\begin{lstlisting}
// 将所有元素按所属集合分组
// 键：根节点编号，值：该组所有元素编号
map<int, vector<int>> groups;
for (int i = 0; i < n; i++) {
    int root = uf.find(i);      // 找到 i 的根节点
    groups[root].push_back(i);  // 将 i 加入对应组
}

// 遍历每个组，处理组内成员
for (auto& entry : groups) {
    int root = entry.first;            // 组标识（根节点）
    vector<int>& members = entry.second;  // 组内所有成员
    // 处理该组 ...
}
\end{lstlisting}

\subsection{应用：朋友与敌人关系}

\begin{lstlisting}
/**
 * 朋友关系 → 合并到同一集合
 * 敌人关系 → 记录到对方的敌人列表中
 * 查询时分别检查是否同集合、是否是敌人
 */

// enemy[i]：存储所有与 i 是敌人关系的节点
vector<vector<int>> enemy(n + 1);
// 并查集初始化
vector<int> parent(n + 1);
for (int i = 1; i <= n; i++) parent[i] = i;

// 处理 m 对关系
for (int i = 0; i < m; i++) {
    int a, b, t;  // a 和 b 的关系类型为 t
    cin >> a >> b >> t;
    if (t == 1) {
        unite(a, b);           // 朋友 → 合并到同一集合
    }
    if (t == -1) {
        enemy[a].push_back(b); // 敌人 → 互相记录
        enemy[b].push_back(a);
    }
}

// 查询 a 和 b 的关系
bool isFriend = same(a, b);     // 同一集合内是朋友
// 在 a 的敌人列表中找 b
bool isEnemy = find(enemy[a].begin(), enemy[a].end(), b)
               != enemy[a].end();

if (isFriend && !isEnemy)    cout << "No problem" << endl;
else if (isFriend && isEnemy) cout << "OK but..." << endl;
else if (!isFriend && isEnemy) cout << "No way" << endl;
else                          cout << "OK" << endl;
\end{lstlisting}

\vfill
\begin{flushright}
\small —— 祝考试顺利！——
\end{flushright}

\end{document}
